\section{结论与建议}

\subsection{主要结论}

\begin{enumerate}[leftmargin=1.9em,itemsep=1pt]
\item \textbf{响应速度是第一变量。} 失联后 95\% 包含域在 1 h 内约 $56\km^2$，
      24 h 增至 $2\,104\km^2$，96 h 达 $27\,652\km^2$。任何装备的覆盖能力都追不上
      不确定性的增长速度。因此“缩短通知–动员–到位时间”应排在任何装备采购之前。
\item \textbf{装备必须匹配目标类别。} 水柱分支占 62\%；只搜海底的侧扫方案 72 h
      成功率仅 38.2\%，而具备水柱探测能力的方案（多波束、AUV）可达 86\%--96\%。
      单套装备中，\textbf{侧扫+多波束（M4，约 182 万美元，89.6\%）}是性价比拐点；
      要求 $\ge95\%$ 则需三件套（M7，约 239 万美元，95.5\%）。
\item \textbf{多资产协同的边际效益有明确顺序。} “加 1 艘同配置船”（+70 万美元）带来
      +7.0 个百分点，是全部选项中性价比最高的一步；“2 船 + 2 AUV”（366 万美元，
      98.9\%）是高保障方案甜点；继续加到 4 船 + 4 AUV（620 万美元）只到 99.9\%，
      边际收益趋零。
\item \textbf{协调规则让“加船”真正有效。} 小时内认领保证同小时重复扫描率为 0；
      重复扫描惩罚 $\beta=1$ 在成功率只降 0.6 个百分点的前提下把跨小时重复率从
      35.6\% 降到 27.5\%。没有协调，多资产会退化为“多船同扫一片”。
\item \textbf{结论可迁移，但不能照搬。} 在加勒比海（小开曼岛西南、开曼海槽边缘）重建
      同一套模型，7 种方案的成功率与地中海相差不超过 2.4 个百分点、排序完全一致；
      但加勒比海表层流速是地中海的 2 倍以上、涡旋更多，在真实预报误差下其实际成功率
      会低于本文结果。强流海域应额外增加冗余并加强现场流场观测。
\end{enumerate}

\subsection{对运营公司的建议}

\begin{enumerate}[leftmargin=1.8em,itemsep=1pt]
\item \textbf{母船按 M4 配置（侧扫 + 多波束 + ROV + 雷达，约 182 万美元）常备}，
      满足 89.6\% 的 72 h 成功率；这也是后续任何扩展的基础；
\item \textbf{与邻近救援船签订 AUV / 多波束快速调用协议}（事故时补成 2 船 + 2 AUV，
      98.9\%），比自购全套更经济；
\item \textbf{把“待命与动员时间 $\le3$ h”写入公司安全规程}，并定期演练；
      这比多买一套装备更有效；
\item \textbf{加勒比等强流海域的运营点}应额外配置 1 艘增援船或 2 台 AUV，
      并建立现场流场观测能力。
\end{enumerate}

\subsection{对监管机构的建议}

\begin{enumerate}[leftmargin=1.8em,itemsep=1pt]
\item \textbf{强制配备 37.5 kHz 声信标（ULB）与应急上浮系统}：这会把“扫海搜索”
      变成“声学定位”，是投入产出比最高的安全措施；
\item \textbf{强制下潜前检查清单与事故报告义务}，并要求运营方持有与事故等级匹配的
      搜救预案（含装备、人员、响应时间承诺）；
\item \textbf{建立与海岸警卫队的联动机制}，把“失联即启动”作为标准程序，
      避免把时间浪费在行政确认上。
\end{enumerate}

\subsection{结语}

本文的模型并不追求预测某一次事故的确切结局，而是回答一个可操作的问题：
\emph{在任何给定时刻，目标可能在哪里、需要多大范围、用哪种装备、多久能找到}。
这条从“位置不确定性”到“时间–成功率–成本”的证据链，在两个差异明显的海域上
都给出了稳定、可解释的结论，可以作为运营与监管决策的依据。

\section*{致希腊政府备忘录}
\addcontentsline{toc}{section}{致希腊政府备忘录}
\markboth{致希腊政府备忘录}{}

\noindent\textbf{主题：小型旅游潜水器失联风险的管理建议}

\vspace{1mm}
\noindent 尊敬的希腊海运与岛屿政策部：

\vspace{1mm}
\noindent 2023 年“泰坦号”事故表明，小型旅游潜水器一旦失联，救援窗口极短。
我们基于伊奥尼亚海与加勒比海两个真实海域的建模，向贵国提出四点建议。

\paragraph{一、承认并量化残余风险。}
我们的模型显示：失联后 24 h，潜水器位于约 $2\,100\km^2$ 范围内的概率为 95\%；
96 h 后扩大到 $27\,652\km^2$——相当于一个 $166\times166$ km 的海域。
“位置已知到 100 m”的乐观假设下尚且如此，实际搜救的难度可想而知。
建议要求运营方在许可证申请中提交“位置不确定性 vs 时间”曲线（即本文 $A_{95}(t)$），
并据此确定最低搜救配置。

\paragraph{二、建立三重防线。}
\begin{enumerate}[leftmargin=1.5em,itemsep=0pt]
\item \textbf{强制 37.5 kHz 声信标（ULB）与应急上浮系统}。声信标可把搜索范围从
      “数百平方公里”压缩到“数平方公里”，是投入产出比最高的单项措施；
\item \textbf{强制遥测与失联告警}：位置、深度、姿态数据实时回传母船，
      失联即自动上报海岸警卫队，杜绝“失联数小时才报警”；
\item \textbf{强制下潜前检查清单与独立安全员制度}，对压载、耐压壳、通信做双人复核。
\end{enumerate}

\paragraph{三、明确搜救能力承诺。}
按本文评估：\textbf{单船“侧扫+多波束”配置}（约 182 万美元）可实现 72 h 内
89.6\% 的发现概率；\textbf{“2 船 + 2 台 AUV”}（约 366 万美元）可达 98.9\%。
建议把“3 h 内首批装备离港、24 h 内形成双船能力”作为运营许可的硬性条件，
并每年演练一次。

\paragraph{四、强化海域协同与审批条款。}
\begin{enumerate}[leftmargin=1.5em,itemsep=0pt]
\item 规定新的旅游潜水器运营点必须提交“海域环境评估”（地形、流场、通信覆盖）；
\item 对强流/涡旋海域（如加勒比海类型的海况）提高装备冗余要求；
\item 建立与海岸警卫队的联合搜救预案，明确指挥权与信息共享机制；
\item 要求事故后 30 天内提交搜救复盘报告，纳入全国数据库。
\end{enumerate}

\vspace{1.5mm}
\noindent 我们的模型与代码完全可复现（仓库附后），可随时按贵国实际海况重算。
愿意为希腊的旅游潜水安全标准制定提供进一步技术支持。

\vspace{3mm}
\noindent 此致

\noindent 敬礼

\vspace{2mm}
\noindent 2024 MCM 参赛队 \quad （供希腊海运与岛屿政策部参考）
